\subsection{向量丛的可微截面的限制}

设 r 是正整数，X 是 \(C^{r}\) 类的局部有限维可微流形，而 E 是一个向量丛（实的或复的），

其基为 X 且是 \(C^{r}\) 类的。对 X 的每个开集 U，用 \(\mathcal{S}^{r}(\mathrm{U};\mathrm{E})\)（在 VAR, R1, 7.4, p. 74 中记作 \(\mathcal{S}_{\mathrm{E}}^{r}(\mathrm{U})\)）表示 E 在 U 上的 \(C^{r}\) 截面所成的向量空间，赋以 \(C^{r}\) 紧收敛拓扑。

引理 4. --- 设 \(\mathcal{U}\) 是 X 的一个开覆盖。映射 \(u: f \mapsto (f|U)_{U \in \mathcal{U}}\)（它从 \(\mathcal{S}^{r}(X;E)\) 到 \(\prod_{U \in \mathcal{U}} \mathcal{S}^{r}(U;E)\) 中）是线性的、单射、连续的且是严格的。它的像是闭的。

映射 \(u\) 是线性的、单射且连续的。依据 TG, IX, p. 43, prop. 1 与 p. 48, cor. 1，X 的每个紧子集都有有限覆盖 \((\mathrm{C}_i)_{i\in \mathrm{I}}\)，使得对每个 \(i\in \mathrm{I}\)，集 \(\mathrm{C}_i\) 是覆盖 \(\mathcal{U}\) 中某个开集的紧子集。由此得出线性映射 \(u\) 是严格的。最后，它的像是闭的，因为它由这样的族 \((f_{\mathrm{U}})_{\mathrm{U}\in \mathcal{U}}\) 组成：\(f_{\mathrm{U}}\) 与 \(f_{\mathrm{V}}\) 在 \(\mathrm{U}\cap \mathrm{V}\) 中重合（对 \(\mathcal{U}\) 中所有的 U 与 V）。

命题 9. --- 空间 \(\mathcal{S}^r(\mathrm{X};\mathrm{E})\) 是完备的。

首先假设存在整数 \(n \geqslant 0\) 与巴拿赫空间 F，使得 X 是 \(R^{n}\) 的开子集，且 E 是以 X 为基、以 F 为纤维的平凡向量丛 \(X \times F\)。在这种情形，拓扑向量空间 \(\mathcal{S}^{r}(X;E)\) 同构于 \(\mathcal{C}^{r}(X;F)\)，结论由 III, p. 34 的引理 3 得出。

在一般情形，设 \(\mathcal{U}\) 是 X 的由坐标域组成的开覆盖，使得对每个 \(\mathrm{U} \in \mathcal{U}\)，E 在 U 上的限制是可平凡化的。由上面的引理，空间 \(\mathcal{S}^r(\mathrm{X};\mathrm{E})\) 同构于线性映射 \(f \mapsto (f|\mathrm{U})_{\mathrm{U} \in \mathcal{U}}\) 的像，而这个像在诸空间 \(\mathcal{S}^r(\mathrm{U};\mathrm{E})\)（其中 U 取遍 \(\mathcal{U}\)）的乘积中是闭的。由前面的情形，每个空间 \(\mathcal{S}^r(\mathrm{U};\mathrm{E})\) 都是完备的，因此它们的乘积是完备的（TG, II, p. 17, prop. 10）。从而 \(\mathcal{S}^r(\mathrm{X};\mathrm{E})\) 是完备的。

注记. --- 设 \(\mathcal{U}\) 是 X 的由坐标域 \(c_{\mathrm{U}} = (\mathrm{U}, \varphi_{\mathrm{U}}, \mathbf{R}^{n_{\mathrm{U}}})\) 组成的开覆盖，使得对每个 \(U \in \mathcal{U}\)，E 在 U 上的限制是 \(F_{U}\) 型可平凡化的，其中 \(F_{U}\) 是巴拿赫空间。对每个 \(U \in \mathcal{U}\)，把 E 在 U 上的截面 f 与映射 \(f_{U}\)（它从 \(\varphi_{\mathrm{U}}(\mathrm{U})\) 到 \(F_{U}\) 中）等同起来。由 prop. 9 的证明可得，空间 \(\mathcal{S}^{r}(\mathrm{X};\mathrm{E})\) 的拓扑由半范数族 \(p_{U,C,\alpha}\) 定义，其中

\[
p _ {\mathrm{U,C}, \alpha} (f) = \sup _ {x \in \mathrm{C}} \| (\partial^ {\alpha} f _ {\mathrm{U}}) (x) \|,
\]

其中 U 取遍 \(\mathcal{U}\)，C 取遍 \(\varphi_{\mathrm{U}}(\mathrm{U})\) 的紧子集全体，而 \(\alpha \in \mathbf{N}^{n_{\mathrm{U}}}\) 取遍满足 \(|\alpha| \leqslant r\) 的多重指标全体。

命题 10. --- 假设向量丛 E 是局部有限秩的。设 s 是满足 \(0 \leqslant s < r\) 的整数，并设 Y 是 X 的相对紧开子集。线性映射 \(f \mapsto f|Y\)（它从 \(\mathcal{S}^{r}(X;E)\) 到 \(\mathcal{S}^{s}(Y;E)\) 中）是紧的。

首先假设存在整数 \(n \geqslant 0\) 与有限维巴拿赫空间 F，使得 X 是 \(R^{n}\) 的开子集，且 E 是以 X 为基、以 F 为纤维的平凡向量丛 \(X \times F\)。在这种情形，拓扑向量空间 \(\mathcal{S}^{r}(X;E)\) 与 \(\mathcal{S}^{s}(Y;E)\) 分别同构于 \(\mathcal{C}^{r}(X;F)\) 与 \(\mathcal{C}^{s}(Y;F)\)，而 prop. 10 是 prop. 8 的推论。

转到一般情形。设 C 是 X 的包含 Y 的紧子集。对 C 的每一点 x，选取 x 在 X 中的一个开邻域 \(U(x)\)，它是一个坐标域，且向量丛 E 在其上是可平凡化的。此外，再选取 x 的一个相对紧开邻域 \(V(x)\)，它包含在 \(U(x)\) 中。由于集 C 是紧的，它被这样一个有限族 \((\mathrm{V}(x_{1}),\ldots,\mathrm{V}(x_{m}))\) 所覆盖。对每个 i，置 \(U_{i}=U(x_{i})\) 与 \(Y_{i}=V(x_{i})\cap Y\)。那么 \(Y_{i}\) 是 \(U_{i}\) 的相对紧开子集，且我们有 \(Y=Y_{1}\cup\cdots\cup Y_{m}\)。线性映射 \(f\mapsto f|U_{i}\)（从 \(\mathcal{S}^{r}(X;E)\) 到 \(\mathcal{S}^{r}(U_{i};E)\) 中）是连续的，而线性映射 \(g\mapsto g|Y_{i}\)（从 \(\mathcal{S}^{r}(U_{i};E)\) 到 \(\mathcal{S}^{s}(Y_{i};E)\) 中）由证明的第一部分是紧的。因此线性映射 \(f\mapsto f|Y_{i}\)（从 \(\mathcal{S}^{r}(X;E)\) 到 \(\mathcal{S}^{s}(Y_{i};E)\) 中）是紧的（III, p. 5 的 prop. 3）。把引理 4（应用于 Y 被诸开集 \(Y_{i}\) 所成的覆盖）与 III, p. 3 的注记 5 和 6 一并考虑，线性映射 \(f\mapsto f|Y\)（从 \(\mathcal{S}^{r}(X;E)\) 到 \(\mathcal{S}^{s}(Y;E)\) 中）是紧的，这就完成了证明。

